\documentclass[aps, superscriptaddress, floatfix, onecolumn]{revtex4-2}
\usepackage{amsmath}
\usepackage{amssymb}
\usepackage{amsthm}
\usepackage{mathtools}
\usepackage{graphicx}
\usepackage{braket}

% For a list of named colors, see https://www.overleaf.com/learn/latex/Using_colors_in_LaTeX
\usepackage{color}
\usepackage[dvipsnames]{xcolor}

%\usepackage{tcolorbox}
%\tcbuselibrary{amsmath}

\newcommand{\result}[1]{\textcolor{blue}{#1}}



\newcommand{\CCQ}{Center for Computational Quantum Physics, Flatiron Institute, 162 5th Avenue, New York, NY 10010, USA}

\newtheorem{remark}{Remark}
\newcommand{\ms}[1]{\textcolor{blue}{[MS: #1]}}
\newtheorem{theorem}{Theorem}

\begin{document}

\title{Gaussian Fermions}

\author{E.\ Miles Stoudenmire}
\affiliation{\CCQ}

\maketitle

\section{Fermionic Vector Spaces}

Fermions are conventionally defined in terms of creation and annihilation operators. Here we will instead start from fermionic vector spaces or state spaces and the properties which make them ``fermionic.'' In the next section, we will see that the more familiar properties of creation and annihilation operators follow from the properties of the states.

The simplest fermionic vector space $V$ has the basis 
\begin{align}
    \mathcal{B}(V) = \{ \ket{0}, \ket{1} \}
\end{align}
which is the basis for a single ``spinless fermion,'' with $\ket{0}$ being the vacuum or empty state and $\ket{1}$ being the filled state or state of one fermion. So far the fermionic nature is not very obvious except  that we can only have at most one fermion. 

Many-body spaces can be formed by taking Cartesian products of spaces (tensor products of basis vectors) so we could consider a space $V_1$ with basis $\{ \ket{0}_1, \ket{1}_1 \}$ and space $V_2$ with basis $\{ \ket{0}_2, \ket{1}_2 \}$, then the space $V_1 \times V_2$ has a basis
\begin{align}
    \mathcal{B}(V_1 \times V_2) = \{ \ket{0}_1 \otimes \ket{0}_2,\  \ket{1}_1 \otimes \ket{0}_2, \  \ket{0}_1 \otimes \ket{1}_2,\  \ket{1}_1  \otimes \ket{1}_2 \} \ .
\end{align}
A generic vector (or ``state'') in this space has the form $\ket{\psi} = a \ket{0}_1 \otimes \ket{0}_2 + b \ket{0}_1 \otimes \ket{1}_2 + c \ket{1}_1 \otimes \ket{0}_2 + d \ket{1}_1 \otimes \ket{1}_2 $.

The defining property of fermions is how one identifies vectors in the space $V_1 \times V_2$  with vectors in the space $V_2 \times V_1$. We can define this identification, or vector space isomorphism, by defining it on the basis vectors. It is given by:
\begin{align}
\ket{0}_1 \otimes \ket{0}_2 & \leftrightarrow {\color{ForestGreen} +}\ket{0}_2 \otimes \ket{0}_1 \nonumber \\
\ket{1}_1 \otimes \ket{0}_2 & \leftrightarrow {\color{ForestGreen} +}\ket{0}_2 \otimes \ket{1}_1 \nonumber \\
\ket{0}_1 \otimes \ket{1}_2 & \leftrightarrow {\color{ForestGreen} +}\ket{1}_2 \otimes \ket{0}_1  \nonumber \\
\ket{1}_1 \otimes \ket{1}_2 & \leftrightarrow {\color{BrickRed} -}\ket{1}_2 \otimes \ket{1}_1  \ \ .
\end{align}
Note the minus sign in the last line. Plus signs are shown for emphasis.

We can summarize the vector space isomorphism as
\begin{align}
     \ket{n_1}_1 \otimes \ket{n_2}_2 \ \leftrightarrow\ (-1)^{n_1 \cdot n_2} \ket{n_2}_2 \otimes \ket{n_1}_1 
\end{align}
where $n_1, n_2 \in \{0,1\}$.
The factor in front of this isomorphism rule means the \emph{ordering} of fermionic basis vectors matters: if two odd basis vectors are permuted, the basis vector changes by a minus sign. 
\footnote{In some mathematical literature, an isomorphism of this type is described as ``having a $\mathbb{Z}_2$ grading'', meaning that one assigns non-trivial properties based on the parity or grading of the basis, or having a ``braided'' or ``graded'' vector space isomorphism.}

Note that tensor products of basis states such as $\ket{1}_i \otimes \ket{1}_j$ are commonly written in the more compact notation $\ket{1}_i \ket{1}_j$ or even $\ket{1_i 1_j}$. Tensor product symbols ``$\otimes$'' are generally dropped when understood from context, and will be dropped frequently in the rest of these notes.


\section{Fermionic Operators}

\subsection{Operator Definitions}

For the simplest case of a ``spinless fermion'' space $V$ with basis $\mathcal{B}(V) = \{ \ket{0}, \ket{1} \}$, we can define the following operators:
\begin{align}
\hat{n}\ &= \ \ket{1} \bra{1} \\ 
\hat{c}\ &= \ \ket{0} \bra{1} \\ 
\hat{c}^\dagger &= \ \ket{1} \bra{0}  \ \ .
\end{align}
Together with the identity operator, this set defines a complete operator basis.

The $\hat{n}$ operator is known as the ``number operator''
since it ``detects'' the filling of a state:
\begin{align}
\hat{n} \ket{0}  &= 0 \cdot \ket{0}  \nonumber \\ 
\hat{n} \ket{1}  &= 1 \cdot \ket{1}  \ \ .
\end{align}

The $\hat{c}$ and $\hat{c}^\dagger$ are known as the ``annihilation'' and ``creation'' operators, respectively. This is simply because they map a filled state to an empty state or vice versa:
\begin{align}
\hat{c} \ket{1}  &= \ket{0}  \\ 
\hat{c}^\dagger \ket{0}  &= \ket{1}  \ \ .
\end{align}

One can of course define tensor products of operators in the usual way, such as $\hat{c}_1 \hat{c_2}$ or $\hat{c}^\dagger_2  \hat{c}_1$.
\footnote{When an operator like $\hat{c}_1$ is written in isolation, it actually stands for $\hat{c}_1 \otimes \hat{I}_2$. Products of operators such as $\hat{c}^\dagger_1 \hat{c}_2$ are actually shorthand for $(\hat{c}^\dagger_1 \otimes \hat{I}_2) (\hat{I}_1 \otimes \hat{c}_2) = \hat{c}^\dagger_1 \otimes \hat{c}_2$.
If the operators act on the same space, they ``stack'', meaning for example that $(\hat{c}^\dagger_1 \otimes \hat{I}_2) (\hat{c}_1 \otimes \hat{I}_2) = (\hat{c}^\dagger_1 \hat{c}_1) \otimes \hat{I}_2$. All of this is easier to see with tensor diagrams!}
These operators act on the spaces $V_1 \times V_2$ or $V_2 \times V_1$, respectively.

\subsection{Commutation Relations}

Now comes a crucial point: the creation and annihilation operators have non-trivial commutation relations, and these are entirely a consequence of the state properties. For operators on state spaces $V_i$ and $V_j$ (with $i \neq j$), the commutation relations are
\begin{align}
\hat{c}_i  \hat{c}_j = -  \hat{c}_j  \hat{c}_i \\ 
\hat{c}^\dagger_i  \hat{c}^\dagger_j = -  \hat{c}^\dagger_j  \hat{c}^\dagger_i \\ 
\hat{c}_i  \hat{c}^\dagger_j = -  \hat{c}^\dagger_j  \hat{c}_i  \label{eqn:ccdag_relation}
\end{align}
The first two of these are commonly summarized as:
\begin{align}
\{ \hat{c}_i, \hat{c}_j \} = 0 \nonumber \\
\{ \hat{c}^\dagger_i, \hat{c}^\dagger_j \} = 0
\end{align}
where $\{ a, b \} = a b + b a$ is the anti-commutator.

The third relation Eq.~(\ref{eqn:ccdag_relation}) has the special case that when $i=j$, $(\hat{c}_i \hat{c}^\dagger_i + \hat{c}^\dagger_i \hat{c}_i) = 1$. Therefore it is commonly summarized as 
\begin{align}
\{ \hat{c}_i, \hat{c}^\dagger_j \} = \delta_{ij} \ \ .
\end{align}

In many presentations, the anticommutation relations of the operators are stated as axioms, but they actually follow from the properties of the state space.
For example, 
\begin{align}
\hat{c}_1  \hat{c}_2   & =  \ket{0}_1 \otimes \bra{1}_1 \otimes \ket{0}_2 \otimes \bra{1}_2 \nonumber \\
& =  \ket{0}_1 \otimes (\bra{1}_1  \otimes \bra{1}_2) \otimes \ket{0}_2 \nonumber \\
& = \ket{0}_1 \otimes ( - \bra{1}_2  \otimes \bra{1}_1) \otimes \ket{0}_2 \nonumber \\
& = - \ket{0}_2 \otimes \bra{1}_2  \otimes \bra{0}_1 \otimes \ket{1}_1 \nonumber \\
& = - \hat{c}_2  \otimes \hat{c}_1 \ \ .
\end{align}

\subsection{Comment on State and Operator Labels}

An important point for what follows is that the subscripts or labels on vectors such as $\ket{0}_i$ or $\ket{1}_j$ or on operators such as $\hat{c}_i$ or $\hat{c}_j$ index over spatial ``sites'' as well as spins, if present. For example, if the fermions are ``spinless fermions'' with sites labeled by integers then the labels are
\begin{align}
i = 1, 2, 3,  4, \ldots
\end{align}
If the fermions are spinless and organized on a two-dimensional lattice then the labels could be
\begin{align}
i = (1,1), (1,2), (1,3), \ldots, (2,1), (2,2), \ldots
\end{align}
Importantly, if the fermions are ``spinful''---electrons being an important example---then the labels are understood to index both sites and spin values
\begin{align}
i = (1,\uparrow), (1,\downarrow), (2,\uparrow), (2, \downarrow), \ldots 
\end{align}
When spin plays an especially important role, such as in computing magnetization or magnetic susceptibility, then it can be helpful to explicitly display spin labels 
such as $\hat{c}_{i,\sigma}$ with $\sigma = \uparrow, \downarrow$.

The key point is that derivations below not explicitly showing spin labels or not explicitly referring to dimensionality are general and apply equally to spinless and spinful fermions in one-, two-, and three-dimensional space etc.\ with an appropriate understanding of the state and operator labels.

\section{Many-Body States and Gaussian States}

Basis states for many-body or many-particle systems can be defined or constructed as
\begin{align}
\ket{1_i 1_j 1_k} = \hat{c}^\dagger_i \hat{c}^\dagger_j \hat{c}^\dagger_k \ket{0} \ .
\end{align}
The most general many-particle state is a sum of these:
\begin{align}
\ket{\Psi} = \sum_{n_1, n_2, \ldots, n_N} \Psi^{n_1 n_2 \cdots n_N} \ket{n_1 n_2 \cdots n_N} \label{eqn:general_state}
\end{align}
where $\Psi^{n_1 n_2 \cdots n_N}$ is an arbitrary tensor of complex amplitudes. The anti-commuting nature of the fermions is automatically satisfied by the properties of the basis, and no additional symmetry property needs to be assumed for the amplitudes $\Psi^{n_1 n_2 \cdots n_N}$. 

A typical additional assumption is that for a closed system the total number of particles $N_p$ is fixed. Mathematically this means $n_1 + n_2 + \ldots + n_N = N_p$ otherwise $\Psi^{n_1 n_2 \cdots n_N} = 0$. (One important exception are states derived from mean-field approximations of superconductors, though one can also formulate descriptions of superconductors which are number-preserving.)

A special subset of many-particle fermionic states are the \emph{Gaussian states}, more commonly known as \emph{Slater determinants}.
The general form of a Gaussian state with four fermions is:
\begin{align}
\ket{\Psi} & = \hat{A}^\dagger \hat{B}^\dagger \hat{C}^\dagger \hat{D}^\dagger \ket{0} \nonumber \\
& = \left( \sum_{i_1=1}^N A^{i_1} \hat{c}^\dagger_{i_1} \right) \left( \sum_{i_2=1}^N B^{i_2} \hat{c}^\dagger_{i_2} \right) \left( \sum_{i_3=1}^N C^{i_3} \hat{c}^\dagger_{i_3} \right) \left( \sum_{i_4=1}^N D^{i_4} \hat{c}^\dagger_{i_4} \right) \ket{0} \ \ \ .
\end{align}
Above we have introduced operators such as $\hat{A} =  \sum_{i_1=1}^N A^{i_1} \hat{c}^\dagger_{i_1}$. One describes the state above by saying that it contains four fermions which are in the ``orbitals'' or ``single-particle wavefunctions'' $\vec{A}, \vec{B}, \vec{C}$ and $\vec{D}$. These orbitals are vectors $\vec{A}= [A^1 \ A^2 \ \ldots A^N]$ which are unit-normalized. The vectors are not required to be orthogonal, though the properties of the state are simpler to understand if they are.

The Gaussian states are quite special and far from being a general fermionic state, even though they are elements of the general vector space of fermionic many-body states. To see that the Gaussian states are not general, let us write a two-particle Gaussian state 
$\ket{AB} = \hat{A}^\dagger \hat{B}^\dagger \ket{0}$ in the general form Eq.~(\ref{eqn:general_state}) for a fermionic state.
We see that for the state $\ket{AB}$, the amplitudes $\Psi^{n_1 n_2 \cdots n_N}$ take the form:
\begin{align}
\Psi^{0_1 0_2 \cdots 1_i \cdots 1_j \cdots 0_N} & = A_i B_j \ \  \text{ for}\ \sum_i n_i = 2\\
\Psi^{n_1 n_2 \cdots n_N} & = 0 \ \ \ \ \  \ \ \text{otherwise for}\ \sum_i n_i \neq 2 \ \ .
\end{align}
So the amplitude tensor $\Psi$ of $\ket{AB}$ is constructed from two vectors (it is rank-2 in a certain sense \footnote{The amplitude tensor for $\ket{AB}$ or a general two-particle Gaussian state is actually a bond-dimension 2 or rank 2 matrix product state}) and is not fully general.

To understand the origin of the name Slater determinant, consider the basis defined as $\ket{i} = \hat{c}^\dagger_i \ket{0}$, the so-called ``real-space basis''.
Expressing states in this basis is known as \emph{first quantization}.
We can express the amplitudes of $\ket{AB}$ in this basis, as $\psi_{AB}(i,j) = \braket{i,j|AB} = (c^\dagger_i c^\dagger_j \ket{0})^\dagger \ket{AB} = \bra{0} c_j c_i \ket{AB}$. One finds that
\begin{align}
\psi_{AB}(i,j) =  A^i B^j - B^i A^j  \label{eq:two_particle_slater}
\end{align}
such that the amplitudes in real space take the form of an antisymmetric function.
It was pointed out by Slater that expressions of the above type can be written as determinants of matrices whose columns are the orbital vectors $\vec{A}, \vec{B}, \ldots$ and whose rows are the spatial positions $i=1,2, \ldots, N$, hence the name Slater determinant. 

The first-quantized basis can be convenient for proving facts such as if $\vec{A} = \vec{B}$, then the state $\ket{AB} = 0$, as is evident from Eq.~(\ref{eq:two_particle_slater}) above. In the more general setting, the requirement that orbitals required be independent is known as the \emph{Pauli exclusion principle}. It is not a separate axiom about fermions, but follows from their vector space properties.


\section{Static Properties of Gaussian States}


In this section we will consider static or time-independent properties of Gaussian states. Here we consider only pure states. 
To define a generic Gaussian state, start by defining a matrix of orbitals $d^{j}\!_m$.
This is assumed to be a unitary matrix $\bar{d}^{m'}\!_j d^j\!_m = \delta^{m'}_m$.
The $m^\text{th}$ column $[\vec{d}_{(m)}]^{j} = d^{j}\!_m$ can be interpreted as the real-space wavefunction
for the single fermion that will be created by the $\hat{d}^\dagger_m$ operator defined below.

Now define the operator $\hat{d}^\dagger_m$ and its conjugate as
\begin{align}
\hat{d}^\dagger_m & = \sum_j d^j\!_m \, \hat{c}^\dagger_j \label{eq:ddagger_transform} \\
\hat{d}^m & = \sum_j \bar{d}^m\!_j \, \hat{c}^j  \ . \label{eq:d_transform}
\end{align}
From the unitary property of $g^j\!_m$ it follows that 
\begin{align}
\hat{c}^\dagger_j & = \sum_m \bar{d}^m\!_j \, \hat{d}^\dagger_m \label{eq:cdagger_transform} \\
\hat{c}^j & = \sum_m d^j\!_m \, \hat{d}^m  \ . \label{eq:c_transform}
\end{align}

The transformation from $\hat{c}$ to $\hat{d}$ operators above is a \emph{canonical transformation}, which just means that the $\hat{d}$ operators obey the same commutation relations and have the same algebraic properties as the $\hat{c}$ operators, such as \mbox{$\{ \hat{d}^m, \hat{d}^\dagger_{m'} \} = \delta^m_{m'}$}. 

Furthermore, although we took the perspective above that the states built from the vacuum $\ket{0}$ by the $\hat{c}^\dagger_j$ operators are our defining, or ``computational'', basis, note the similarity of the transformations (\ref{eq:ddagger_transform}), (\ref{eq:d_transform}) to the transformations (\ref{eq:cdagger_transform}), (\ref{eq:c_transform}). One can just as equally view the states built up from the $\hat{d}^\dagger_m$ operators as the defining basis and conclude that properties like expected values of $\hat{c}$ or $\hat{c}^\dagger$ operators in such states will be equivalent up to replacing $d^j\!_m \leftrightarrow \bar{d}^m\!_j$.

\subsection{Generic Gaussian State}

The most generic form of a pure Gaussian state or Slater determinant is given by
\begin{align}
\ket{\mathbf{m}}_d = \ket{m_1 m_2 \ldots m_{N}}_d & = \hat{d}^\dagger_{m_1} \hat{d}^\dagger_{m_2} \cdots \hat{d}^\dagger_{m_{N}} \ket{0} \\
& = \left[ \prod_{m=1} (\hat{d}^\dagger_{m})^{\eta_m} \right] \ket{0}
\end{align}
for some set of orbitals $d^j\!_m$ where $\eta_m$ is the \emph{occupancy} of state $m$ in the $d$ basis. The occupancy is defined here as $\eta_m = 1$ if $m \in \{m_1, m_2, \ldots, m_n\}$ and $\eta_m = 0$ otherwise. The total number of fermions is $N$.

\subsection{Correlation matrix}

A central quantity for Gaussian states is the correlation matrix $C_i\!^j = \bra{\mathbf{m}}_d  \hat{c}^\dagger_i \hat{c}^j \ket{\mathbf{m}}_d $. We will see  that knowledge of the correlation matrix is sufficient to compute expected values of operators, the action
of operators on Gaussian states, and dynamical properties such as single-particle Green's functions \footnote{This observation is closely analogous to the fact that  correlation functions of multi-variate Gaussian functions $\exp(- x^\dagger \Sigma^{-1} x)$ are uniquely determined by their covariance matrix $\Sigma$.}.

The correlation matrix of a Gaussian state $\ket{\mathbf{m}}_d$ is given by
\begin{align}
\boxed{C_i\!^j =  \bra{\mathbf{m}}_d  \hat{c}^\dagger_i \hat{c}^j \ket{\mathbf{m}}_d = \sum_m \bar{d}_i^m \eta_m d^j_m} \ \ . \label{eq:correlation_matrix}
\end{align}
where recall that the $d^j_m$ are the orbitals and $\eta_m$ the occupancies defining the state.

To derive this result, note that for a Gaussian state $\ket{\mathbf{n}}_c  = \left[ \prod_{n=1} (\hat{c}^\dagger_{n})^{\gamma_n} \right] \ket{0}$
the correlation matrix is \mbox{$\bra{\mathbf{n}}_c \hat{c}^\dagger_i \hat{c}^j  \ket{\mathbf{n}}_c = \delta_i^j \gamma_i$}, i.e.\ it is a diagonal matrix of the occupations.
This result follows from the fact
that $\hat{c}^j  \ket{\mathbf{n}}_c = \gamma_j \ket{\mathbf{n'}}_c$ where $\mathbf{n'}$ is the set $\mathbf{n}$ with the element $j$ removed if present.
Then recalling that the transformation from the $\hat{c}$ operators to the $\hat{d}$ operators is a canonical transformation, a similar
result holds in the $d$ basis, namely $\bra{\mathbf{m}}_d \hat{d}^\dagger_m \hat{d}^{m'}  \ket{\mathbf{m}}_d = \delta_m^{m'} \eta_m$.
Finally to derive Eq.~(\ref{eq:correlation_matrix}) use (\ref{eq:cdagger_transform}) and (\ref{eq:c_transform}) to write 
\begin{align}
C_i\!^j = \bra{\mathbf{m}}_d  \hat{c}^\dagger_i \hat{c}^j \ket{\mathbf{m}}_d = \sum_{m,m'} \bar{d}^m\!_i d^j\!_{m'}  \bra{\mathbf{m}}_d  \hat{d}^\dagger_m \hat{d}^{m'} \ket{\mathbf{m}}_d  = \sum_m \bar{d}_i^m \eta_m d^j_m \ \ . \qed
\end{align}

For a normalized state $\ket{\mathbf{m}}_d$, note that the trace of the correlation matrix gives the number of fermions in the state 
\begin{align}
N \stackrel{\text{def}}{=} \sum_j \bra{\mathbf{m}}_d \hat{n}_j \ket{\mathbf{m}}_d = \sum_j \bra{\mathbf{m}}_d \hat{c}^\dagger_j \hat{c}^j \ket{\mathbf{m}}_d = \sum_j C_j\!^j = \sum_m \eta_m \ .
\end{align}
If the state is pure, the occupations $\eta_m$ are either 0 or 1, so that the correlation matrix is a projector $C^2 = C$. 

\subsection{Expected Value of Quadratic Operators}

Given a quadratic operator $\hat{\mathcal{O}} = \sum_{i,j} O^i\!_j \hat{c}^\dagger_i c^j$, we can use the correlation matrix result above to derive the expected value 
$\bra{\mathbf{m}}_d \hat{\mathcal{O}} \ket{\mathbf{m}}_d$. It follows by linearity that
\begin{align}
\boxed{\bra{\mathbf{m}}_d \hat{\mathcal{O}} \ket{\mathbf{m}}_d = \sum_{i,j} O^i\!_j \bra{\mathbf{m}}_d  \hat{c}^\dagger_i c^j\ket{\mathbf{m}}_d = \sum_{i,j} O^i\!_j C_i\!^j = \text{Tr}[O C] } \ \ .
\end{align}
An important example of a quadratic operator is a non-interacting Hamiltonian, which in the particle-conserving case has the generic form $\hat{H} =  \sum_{i,j} h^i\!_j \hat{c}^\dagger_i \hat{c}^j$. So from the result above we see that the energy expectation value of a Gaussian state depends only on its correlation matrix.

\subsection{Action of Operators on Gaussian States}

Given an arbitrary creation operator $\hat{v}^\dagger = v^j \hat{c}^\dagger_i$, where $v^j$ is some vector of coefficients, consider its action on a Gaussian state $\ket{\phi}$, resulting in the state $\ket{\tilde{\psi}} = \hat{v}^\dagger \ket{\phi}$.
The state $\ket{\tilde{\psi}}$ is a Gaussian state and its correlation matrix $( C_{\tilde{\psi}})_i\!^j = \bra{\tilde{\psi}} \hat{c}^\dagger_i \hat{c}^j \ket{\tilde{\psi}}$  is given by 
\begin{align}
\boxed{C_{\tilde{\psi}} = ||v_0||^2 C_\phi + v_0 v_0^\dagger}
\end{align}
where $C_\phi$ is the correlation matrix of $\ket{\phi}$ and $v_0 = (1-C_\phi)v$ is the projection of $v$ onto the unoccupied subspace of $C_\phi$.


Analogously, given an arbitrary annihilation operator $\hat{w} = \bar{w}_j \hat{c}^i$, where $w^j$ is some vector of coefficients, consider its action on a Gaussian state $\ket{\phi}$, resulting in the state $\ket{\tilde{\eta}} = \hat{w} \ket{\phi}$.
The state $\ket{\tilde{\eta}}$ is a Gaussian state and its correlation matrix $( C_{\tilde{\eta}})_i\!^j = \bra{\tilde{\eta}} \hat{c}^\dagger_i \hat{c}^j \ket{\tilde{\eta}}$  is given by 
\begin{align}
\boxed{C_{\tilde{\eta}} = ||w_1||^2 C_\phi - w_1 w_1^\dagger}
\end{align}
where $C_\phi$ is the correlation matrix of $\ket{\phi}$ and $w_1 = C_\phi w$ is the projection of $w$ onto the occupied subspace of $C_\phi$.

Note that the states $\ket{\tilde{\psi}}$ or $\ket{\tilde{\eta}}$ are not normalized. This results in the trace of their correlation matrices $C_{\tilde{\psi}}$ or $C_{\tilde{\eta}}$ no longer having a trace equal to the correct particle number. (For $\ket{\tilde{\psi}}$ the correct particle number should be $N_\phi + 1$ and for $\ket{\tilde{\eta}}$ the correct particle number should be $N_\phi-1$.) If desired, the correct normalization can be restored by multiplying by the appropriate factor.


\subsection{Inner Product of Gaussian States}

Consider two Gaussian pure states $\ket{\mathbf{m}}_d$ and $\ket{\mathbf{n}}_f$ given by:
\begin{align}
\ket{\mathbf{m}}_d & = \ket{m_1 m_2 \ldots m_{N}}_d = \hat{d}^\dagger_{m_1} \hat{d}^\dagger_{m_2} \cdots \hat{d}^\dagger_{m_{N}} \ket{0} \\
\ket{\mathbf{n}}_f & = \ket{n_1 n_2 \ldots n_{N}}_f = \hat{f}^\dagger_{n_1} \hat{f}^\dagger_{n_2} \cdots \hat{f}^\dagger_{n_{N}} \ket{0} \\
\end{align}
where the number of particles $N$ is taken to be the same. (Otherwise the inner product of these states is zero.)


\section{Dynamics of Gaussian States}

An appealing aspect of Gaussian states is that their dynamical or time-dependent properties can be completely understood through simple and efficiently- computable matrix expressions. In this section we will consider dynamics governed by a quadratic Hamiltonian of the form 
\begin{align}
\hat{H} = \sum_{ij} h^i\!_j \hat{c}^\dagger_i \hat{c}^j \ \ .
\end{align} 
We will construct its eigenoperators and eigenstates then discuss how generic Gaussian states evolve under this Hamiltonian and show expressions for the Green's functions (time-dependent correlation functions) of these states.

\subsection{Eigenoperators}

To construct eigenoperators, begin by diagonalizing the Hamiltonian ``hopping'' or coefficient matrix $h^i\!_j$ as
\begin{align}
h^i\!_j = \phi^i\!_n \epsilon_n \bar{\phi}_j\!^n \ \ .
\end{align}
The resulting ``single-particle eigenstates'' $\phi^i\!_n$ will play a crucial role in what follows. These can be viewed as spatial single-particle wavefunctions or orbitals labeled by $n$.

Next construct eigenoperators defined as
\begin{align}
\hat{\phi}^\dagger_n & = \sum_j \phi^j\!_n \hat{c}^\dagger_j \label{eq:phidag_op} \\ 
\hat{\phi}^n & = \sum_j \bar{\phi}_j\!^n \hat{c}^j \label{eq:phi_op}  \ \ .
\end{align}
from which it follows by unitarity of $\phi^j_n$ that 
\begin{align}
\hat{c}^\dagger_j & = \sum_n \bar{\phi}_j\!^n \hat{\phi}^\dagger_n  \label{eq:cdag_phi_transform} \\ 
\hat{c}^j & = \sum_n \phi^j\!_n \hat{\phi}^n \ \ . \label{eq:c_phi_transform}
\end{align}

The reason we are calling these eigenoperators is that they only transform by a phase when evolving in time in the Heisenberg picture
\begin{align}
\boxed{\hat{\phi}^n(t) =  e^{-i \epsilon_n t}\ \hat{\phi}^n(0)} \label{eq:phi_t_evol}
\end{align}
\begin{align}
\boxed{\hat{\phi}^\dagger_n(t) =  e^{i \epsilon_n t}\ \hat{\phi}^\dagger_n(0)} \label{eq:phidag_t_evol} \ \ .
\end{align}
To show this result, use the fact that the $\hat{\phi}^\dagger_n$ and $\hat{\phi}^n$ obey the canonical fermionic anticommutation relations and use this to show $[\hat{\phi}^n, \hat{H}] = \epsilon_n \hat{\phi}^n$ and $[\hat{\phi}^\dagger_n, \hat{H}] = -\epsilon_n \hat{\phi}^\dagger_n$. These commutation relations hold for $\hat{\phi}^n(t)$ at all times. Then the properties (\ref{eq:phi_t_evol}) and (\ref{eq:phidag_t_evol}) follow from the Heisenberg equation $i \frac{d}{dt} \hat{\phi}^n(t) = [\hat{\phi}^n(t), \hat{H}]$ and similarly for $\hat{\phi}^\dagger_n(t)$.

\subsection{Eigenstates}

Having found eigenoperators, we can construct many-body, Gaussian eigenstates. These are generically of the form
\begin{align}
\ket{\mathbf{n}}_\phi = \prod_{n} \left(\hat{\phi}^\dagger_n \right)^{\nu_n} \ket{0}  
\end{align}
where $\nu_n = 1$ if $n \in \{n_1, n_2, \ldots, n_N \}$ and $\nu_n = 0$ otherwise. Such states, known as \emph{Fock states}, 
evolve only by a phase under Schr\"odinger time evolution
\begin{align}
\boxed{e^{-i \hat{H} t} \ket{\mathbf{n}}_\phi  = e^{-i E_\mathbf{n} t}  \ket{\mathbf{n}}_\phi }
\end{align}
where 
\begin{align}
\boxed{E_\mathbf{n}  = \epsilon_{n_1} + \epsilon_{n_2} + \ldots + \epsilon_{n_N}} \label{eq:fock_energy}
\end{align}
 is the energy of the state.
 
To derive this result, we can use the Heisenberg picture to write
\begin{align}
e^{-i \hat{H} t} \ket{\mathbf{n}}_\phi  & = e^{-i \hat{H} t} \hat{\phi}^\dagger_{n_1} \hat{\phi}^\dagger_{n_2} \cdots \hat{\phi}^\dagger_{n_N} \ket{0} \\
& =  \left( e^{-i \hat{H} t} \hat{\phi}^\dagger_{n_1}  e^{i \hat{H} t} \right) \left(e^{-i \hat{H} t} \hat{\phi}^\dagger_{n_2}  e^{i \hat{H} t} \right) \cdots  \left( e^{-i \hat{H} t} \hat{\phi}^\dagger_{n_N}  e^{i \hat{H} t}\right) \ket{0}  \\ 
& = e^{-i (\epsilon_{n_1} + \epsilon_{n_2} + \ldots + \epsilon_{n_N}) t} \ket{\mathbf{n}}_\phi \label{eq:eigenstate_t_evolution}
\end{align}
where we have used $e^{-i \hat{H} t} \hat{\phi}^\dagger_{n}  e^{i \hat{H} t} = \hat{\phi}_n(-t) =  e^{-i \epsilon_n t} \hat{\phi}_n(0)$ and $e^{i \hat{H} t} \ket{0} = \ket{0}$.
 
Equation~(\ref{eq:fock_energy}) is a crucial result, and justifies the level-filling picture used in band theory. It says we can understand the full, many-body spectrum of a Gaussian fermionic Hamiltonian just by enumerating the single-particle states $\phi_n^j$ for $n=1,...,V$ where $V$ is the number of sites (volume) of the system. Then the $2^{N_S}$ different many-body eigenstates are just the exponentially many different ways of ``filling'' these $N_S$ single particle states or levels.


\subsection{Zero-Temperature Green's Functions}

With these results in hand, it is straightforward to find expressions for the zero-temperature or ground state Green's functions. 
Their definitions are general for any many-body system while the expressions for them in terms of the $\phi^j_n$ matrix is specific to Gaussian states.

The lesser Green's function is given by
\begin{align}
\boxed{G^{<\, i}_j(t)= i \bra{\Phi_0} \hat{c}^\dagger_j \hat{c}^i(t) \ket{\Phi_0}  = i \sum_n \nu_n  \phi^i_n e^{-i \epsilon_n t}  \bar{\phi}_j^n} \ \ .
\end{align}
The greater Green's function is given by
\begin{align}
\boxed{G^{>\, i}_j(t)= -i \bra{\Phi_0} \hat{c}^i(t) \hat{c}^\dagger_j \ket{\Phi_0} = -i \sum_n (1-\nu_n) \phi^i_n e^{-i \epsilon_n t} \bar{\phi}_j^n} \ \ .
\end{align}
Finally the retarded Green's function is given by
\begin{align}
\boxed{G^{R\, i}_j(t) = \theta(t)(G^{>\, i}_j(t) - G^{<\, i}_j(t)) = -i \theta(t) \sum_n \phi^i_n e^{-i \epsilon_n t} \bar{\phi}_j^n  = -i \theta(t) e^{-i \mathbf{h} t} } \ \ .  \label{eq:GR}
\end{align}
where $\theta(t) = 1$ for $t > 0$ and $\theta(t) = 0$ otherwise.
In all the results above, the $\nu_n$ are the occupation numbers of the ground state of $N$ fermions, which means they are $\nu_n = 1$ for the $N$ lowest-energy single-particle eigenstates and $\nu_n = 0$ otherwise.

These results can all be derived using Eqs.~(\ref{eq:phidag_op}) and (\ref{eq:phi_op}) to transform the $\hat{c}$ and $\hat{c}^\dagger$ into linear combinations of the $\hat{\phi}$, $\hat{\phi}^\dagger$ operators. Then the results above follow from the time evolution of the $\hat{\phi}$ operators Eqs.~(\ref{eq:phi_t_evol}) and (\ref{eq:phidag_t_evol}).

The factor of $-i$ in the definition of $G^R$ helps its Fourier transform have a particularly simple form. The Fourier transform of $f(t) = \theta(t) e^{-i \epsilon t}$ is $\tilde{f}(w) = i/(\omega - \epsilon + i \alpha)$ with $\alpha > 0$ a small broadening factor. Using this fact in Eq.~(\ref{eq:GR}) gives the expression for $G^R$ in frequency space
\begin{align}
\boxed{ \tilde{G}^{R\, i}_j(\omega)  = \sum_n \frac{\phi^i_n \bar{\phi}_j^n}{\omega - \epsilon_n + i \alpha} = \frac{1}{\omega - \mathbf{h} + i \alpha} } \ \ .
\end{align}
This form is known as the Lehmann representation. In the last expression above, the terms in the denominator are matrices and it is understood that the $\omega$ and $i \alpha$ factors are multiplied by identity matrices.

\subsection{Time Evolution of a Generic Gaussian State}

Consider now a generic Gaussian state
\begin{align}
\ket{\mathbf{m}}_d = \ket{m_1 m_2 \ldots m_{N}}_d & = \hat{d}^\dagger_{m_1} \hat{d}^\dagger_{m_2} \cdots \hat{d}^\dagger_{m_{N}} \ket{0} \\
& = \left[ \prod_{m=1} (\hat{d}^\dagger_{m})^{\eta_m} \right] \ket{0} \ \ .
\end{align}
where the $\hat{d}^\dagger_m$ operators are defined as
\begin{align}
\hat{d}^\dagger_m = \sum_j d^j\!_m \hat{c}^\dagger_j \ \ . \label{eq:ddag_c_transform}
\end{align}
Define the \emph{propagator} as $\boxed{g(t) = e^{- i \mathbf{h} t}}$ in other words $\boxed{g^i\!_j(t) = \sum_n \phi^i\!_n e^{-i \epsilon_n t} \bar{\phi}^n\!_j}$.

Using the fact that $\hat{d}^\dagger_m = \sum_j d^j_m \hat{c}^\dagger_j = \sum_{j, n} (\bar{\phi}^n_j d^j_m ) \hat{\phi}^\dagger_n$, where we have used equations (\ref{eq:ddag_c_transform}) and (\ref{eq:cdag_phi_transform}), it follows that 
\begin{align}
\ket{\mathbf{m}}_d & = \sum_{j_1 j_2 \cdots}   \sum_{n_1 n_2 \cdots}  ( \bar{\phi}^{n_1}_{j_1} d^{j_1}_{m_1} ) (\bar{\phi}^{n_2}_{j_2} d^{j_2}_{m_2} ) \cdots (\bar{\phi}^{n_N}_{j_N} d^{j_N}_{m_N} ) \ \hat{\phi}^\dagger_{n_1} \hat{\phi}^\dagger_{n_2}  \cdots \hat{\phi}^\dagger_{n_N} \ket{0} 
\end{align}
which implies, using Eq.~(\ref{eq:eigenstate_t_evolution}) that
\begin{align}
e^{-i \hat{H} t} \ket{\mathbf{m}}_d & = \sum_{j_1 j_2 \cdots}   \sum_{n_1 n_2 \cdots}  ( \bar{\phi}^{n_1}_{j_1} d^{j_1}_{m_1} ) (\bar{\phi}^{n_2}_{j_2} d^{j_2}_{m_2} ) \cdots (\bar{\phi}^{n_N}_{j_N} d^{j_N}_{m_N} ) \ e^{-i (\epsilon_{n_1} + \epsilon_{n_2} + \ldots) t}\  \hat{\phi}^\dagger_{n_1} \hat{\phi}^\dagger_{n_2}  \cdots \hat{\phi}^\dagger_{n_N} \ket{0}  \\ 
& = \sum_{j_1 j_2 \cdots}   \sum_{n_1 n_2 \cdots}  (\phi^{j'_1}_{n_1} e^{-i \epsilon_{n_1} t}   \bar{\phi}^{n_1}_{j_1} d^{j_1}_{m_1}) (\phi^{j'_2}_{n_2} e^{-i \epsilon_{n_2} t}  \bar{\phi}^{n_2}_{j_2}  d^{j_2}_{m_2}) \cdots (\phi^{j'_N}_{n_N} e^{-i \epsilon_{n_N} t}  \bar{\phi}^{n_N}_{j_N}  d^{j_N}_{m_N})  \ \hat{c}^\dagger_{n_1} \hat{c}^\dagger_{n_2}  \cdots \hat{c}^\dagger_{n_N} \ket{0}   \\
& = \sum_{j_1 j_2 \cdots}   \sum_{n_1 n_2 \cdots}  (g^{j'_1}\!_{j_1}(t) d^{j_1}_{m_1}) (g^{j'_2}\!_{j_2}(t) d^{j_2}_{m_2}) \cdots (g^{j'_N}\!_{j_N}(t) d^{j_N}_{m_N})  \ \hat{c}^\dagger_{n_1} \hat{c}^\dagger_{n_2}  \cdots \hat{c}^\dagger_{n_N} \ket{0}   \\
& = \hat{d}^\dagger_{m_1}(t) \hat{d}^\dagger_{m_2}(t) \cdots \hat{d}^\dagger_{m_N}(t)  \ket{0} \ \ .
\end{align}
where we have defined
\begin{align}
\boxed{\hat{d}^\dagger_m(t) = \sum_{j'} d^{j'}\!_m(t) \hat{c}^\dagger_{j'} = \sum_{j' j} \left(g^{j'}\!_{j}(t) d^j\!_m \right) \hat{c}^\dagger_{j'}} \ \ . 
\end{align}
We therefore see that generic orbitals, defined by a matrix of orbitals $d^j\!_m$, evolve in time by acting with the propagator matrix $g^{j'}\!_j(t)$ onto
the orbital matrix $d^j\!_m$.

As a special case of this result, it is interesting to observe that if we choose the orbitals to be the eigenorbitals $d^j\!_m = \phi^j\!_m$ then we find 
$\phi^j\!_m(t) = g^{j'}\!_j(t) \phi^j\!_m = e^{-i \epsilon_m t} \phi^j\!_m$.

\subsection{Time Evolution of a Generic Correlation Matrix}

As we saw above, for a Gaussian state $\ket{\mathbf{m}}_d$ with correlation matrix $C_i\!^j = \bra{\mathbf{m}}_d \hat{c}^\dagger_i \hat{c}^j \ket{\mathbf{m}}_d = \sum_m \bar{d}_j\!^m \eta_m d^j\!_m$, the orbitals evolve, in the Schr\"odinger picture, as $d^i\!_m(t) = g^i\!_j(t) d^j\!_m$. 
So we find that 
\begin{align}
\boxed{C_i\!^j(t) =  \bar{g}^{i'}_i(t)\, C_{i'}\!^{j'}\, g^j_{j'}(t) } \ \ .
\end{align}


\section{Thermal Gaussian States}

The Gaussian formalism extends straightforwardly to finite temperature, also offering closed-form results for all properties of interest. Given a Gaussian Hamiltonian $\hat{H}= \sum_{i,j} h^i_j \hat{c}^\dagger_i \hat{c}^j$, the finite-temperature (mixed) state in the canonical ensemble formalism at temperature $T$ or coolness $\beta = 1/T$ has the form
\begin{align}
\hat{\rho} = \frac{1}{Z}\ e^{-\beta \hat{H}} \ \ .
\end{align}
Given a spectral decomposition of the ``hopping matrix'' $h^i_j = \sum_n \phi^i_n \epsilon_n \bar{\phi}^n_j$ and defining $\hat{\phi}^\dagger_n = \sum_j \phi^j_n \hat{c}^\dagger_j$, we can then write 
\begin{align}
\hat{\rho} = \frac{1}{Z}\   e^{-\beta \sum_n  \epsilon_n \hat{\phi}^\dagger_n \hat{\phi}^n} \ \ . 
\end{align}

\subsection{Thermal Correlation Matrix}

The correlation matrix of a thermal Gaussian state is 
\begin{align}
    C_i\,^j = \text{Tr}[ \hat{c}^\dagger_i \hat{c}^j ] = \left[ \frac{1}{1+e^{\beta \mathbf{h}}} \right]_i^j  =  \sum_n \bar{\phi}^n_i \, f_\beta(\epsilon_n) \, \phi^j_n 
\end{align}
where $f_\beta(\epsilon)$ is the \emph{Fermi function}
\begin{align}
f_\beta(\epsilon) = \frac{1}{1+e^{\beta \epsilon}} \ \ .
\end{align}
Notably, for low temperatures $\beta \rightarrow \infty$, $f_\beta(\epsilon \leq 0) \rightarrow 1$ and $f_\beta(\epsilon > 0) \rightarrow 0$.


\section{Entanglement and Schmidt Decomposition of Gaussian States}

Consider a pure Gaussian state $\ket{\psi}$ with correlation matrix $C_i\,^j = \bra{\psi} \hat{c}^\dagger_i \hat{c}^j \ket{\psi}$. 
Our goal is to obtain a Schmidt decomposition across two disjoint regions $A$ and $B$ meaning 
\begin{align}
\ket{\psi} = \sum_\mu w_\mu \ket{\mu}_A \ket{\mu}_B \label{eq:schmidt}
\end{align}
where $\ket{\mu}_A$ is an orthonormal basis of $A$ and  $\ket{\mu}_B$ an orthonormal basis of $B$. 
To find this decomposition, it will turn out to be sufficient to diagonalize the block of $C$ restricted to $A$, called $C_{AA}$. Below we will find an
expression for the Schmidt values or weights $w_\mu$ in terms of the eigenvalues of $C_{AA}$.

\subsection{Properties of Correlation Matrix Blocks}

First we will establish some useful relationships between blocks of a correlation matrix. The only assumptions needed will be that the correlation matrix of
a pure Gaussian state is a projector and is Hermitian:
\begin{align}
C^2 & = C \label{eq:projector} \\
C^\dagger & = C \ \ . \label{eq:hermitian}
\end{align}

Consider the blocks of $C$ formed by restricting its rows and columns to regions $A$ or $B$ (permuting the rows and columns if needed):
\begin{align}
C = \begin{bmatrix}
C_{AA} & C_{AB} \\ 
C_{BA} & C_{BB} 
\end{bmatrix}
\end{align}
Because $C^\dagger = C$, it follows that $C^\dagger_{AB} = C_{BA}$ and that $C_{AA}$ and $C_{BB}$ are Hermitian.
Then the property $C^2=C$ implies the following relations
\begin{align}
C_{AB} C_{BA} & = C_{AA} (1- C_{AA})  \label{eq:proj1} \\
C_{BA} C_{AB} & = C_{BB} (1 - C_{BB}) \\
C_{AA} C_{AB} & = C_{AB} (1 - C_{BB}) \\
C_{BB} C_{BA} & = C_{BA} (1 - C_{AA})  \ \ .
\end{align}

\subsection{Block Decomposition}

Without loss of generality, take the linear size $m$ of the space $B$ to be greater than or equal to linear size $n$ of $A$, so $m \leq n$.
Then $C_{AA}$ is an $n \times n$ matrix, $C_{BB}$ is $m \times m$, and $C_{AB}$ is $n \times m$.

To obtain the Schmidt decomposition Eq.~(\ref{eq:schmidt}), first compute the singular value decomposition of $C_{AB}$
\begin{align}
C_{AB} = A \Lambda B^\dagger \ \ .
\end{align}
where $\Lambda = \text{diag}(\lambda_1, \lambda_2, \ldots)$ is a square, real, non-negative, diagonal matrix of linear dimension $r$ where
$r = \min(m,n)$.
The matrices  $A$ and $B$ are $n \times r$ and $m \times r$, respectively, and obey
$B^\dagger B = 1$ and $A^\dagger A = 1$.

Note that $C_{AB} C_{BA} = A \Lambda^2 A^\dagger$. By using Eq.~(\ref{eq:proj1}) we can obtain the relation:
\begin{align}
C_{AA} (1-C_{AA}) = A \Lambda^2 A^\dagger \ \ . \label{eq:svd_relation}
\end{align}
Thus $C_{AA}$ is also diagonalized by the matrix $A$
\begin{align}
C_{AA} = A \Omega A^\dagger
\end{align}
where $\Omega$ is a diagonal matrix $\Omega = \text{diag}(\nu_1, \nu_2, \ldots)$ of ``occupations'' ($r$ of these). 
Equation~(\ref{eq:svd_relation}) implies the eigenvalues $\nu_j$ of $C_{AA}$ are related to the 
singular values $\lambda_j$ of $C_{AB}$ through 
\begin{align}
\lambda_j = \sqrt{\nu_j (1-\nu_j)} \ \ .
\end{align}

Since $C$ is a projector, one can show that the eigenvalues of diagonal blocks such as $C_{AA}$ lie between 0 and 1. 
Each $\nu_j$ is within the range $0 \leq \nu_j \leq 1$ so the singular values of $C_{AB}$ are in the range $0 \leq \lambda_j \leq 1/2$.

To complete the analysis of the blocks, we likewise find that
\begin{align}
C_{BB} (1-C_{BB}) = B \Lambda^2 B^\dagger
\end{align}
To make the above equation consistent with the decompositions of $C_{AB}$ and $C_{AA}$ above, we can conclude that $C_{BB}$, which we took to
be as large or larger than $C_{AA}$ has two kinds of eigenvalues:
\begin{enumerate}
\item Eigenvalues $\tilde{\nu}_j = (1-\nu_j)$ which fall between $0 \leq \tilde{\nu}_j < 1$, where recall $\nu_j$ is an eigenvalue of $C_{AA}$.
\item Eigenvalues $\tilde{\nu}_j = 1$.
\end{enumerate}
One can understand the second case of eigenvectors of $C_{BB}$ which have eigenvalue $\tilde{\nu}_j = 1$ as spanning the space not captured by 
the column space of the rectangular matrix $B$.
These $\tilde{\nu}_j = 1$ orbitals have spatial support which is entirely inside of region $B$.

\subsection{Construction of Filled and Empty Orbitals}

From the above analysis we can construct the Schmidt states by defining operators
\begin{align}
\hat{\psi}^\dagger_j = \sum_{j'} \psi^{(1) j'}_j \hat{c}^\dagger_{j'} \ \ .
\end{align}
As will be shown below, these operators create ``occupied natural orbitals''---orbitals which are
eigenvectors of the correlation matrix $C$ with eigenvalue 1---which implies that
the state can be written as:
\begin{align}
\ket{\psi} = \prod_{j} \hat{\psi}^\dagger_j \ket{0}
\end{align}
Using the above analysis of the correlation matrix and defining columns of the $A$ and $B$ matrices as
\begin{align}
(\vec{a}_j)^{i} = A^i\,_j  \\
(\vec{b}_j)^{i} = B^i\,_j  \\
\end{align}
the occupied orbitals $(\vec{\psi}^{(1)}_j)^{j'} = \psi^{(1) j'}_j$ are 
\begin{align}
\vec{\psi}^{(1)}_j = \begin{bmatrix}
\sqrt{\nu_j} & \vec{a}_j \\
\sqrt{1-\nu_j} &  \vec{b}_j \\
\end{bmatrix} \label{eq:occupied_orbital}
\end{align}
For completeness, one can also find the unoccupied orbitals to be 
\begin{align}
\vec{\psi}^{(0)}_j = \begin{bmatrix}
\sqrt{1-\nu_j} & \vec{a}_j \\
- \sqrt{\nu_j} &  \vec{b}_j \\
\end{bmatrix}
\end{align}
Using (\ref{eq:occupied_orbital}), the occupied orbital creation operators can be written as
\begin{align}
\hat{\psi}^\dagger_j & = \sqrt{\nu_j} \ \hat{a}^\dagger_j + \sqrt{1-\nu_j} \ \hat{b}^\dagger_{j} \\
\end{align}
where
\begin{align}
\hat{a}^\dagger_j & = \sum_{i \in A} (\vec{a}_j)^i \hat{c}^\dagger_i \\
\hat{b}^\dagger_j & = \sum_{i \in B} (\vec{b}_j)^i \hat{c}^\dagger_i
\end{align}
are creation operators supported in regions $A$ and $B$ respectively.

\end{document}
